\begin{ouverture}[geometrie]{17}{Solides, disques et volumes}{Partie 3 · Espace et géométrie}
{Perspectives · Patrons · Prismes · Cylindres · Capacités}
\objectif{Relier perspectives, vues et patrons.}
\objectif{Calculer les volumes usuels et l’aire du disque.}
\objectif{Convertir volumes et capacités.}
\prerequis{Aires, carrés et cubes.}
\end{ouverture}
\pagestyle{otMain}\essentiel\fontsize{11}{15}\selectfont
\phantomsection\label{cours-17-1}
\essSection{1}{Représenter l’espace}
\begin{definition}Un pavé droit a six faces rectangulaires ; un cube a six faces carrées. Un prisme droit a deux bases polygonales identiques et parallèles, et des faces latérales rectangulaires. Un cylindre de révolution a deux bases en forme de disques et une surface latérale courbe.\end{definition}
\label{fig-cours-17-1-cours-solides}\figureGeo{cours-solides}

\begin{savoirfaire}[Lire une représentation]En perspective cavalière, les arêtes parallèles restent parallèles sur le dessin. Les arêtes cachées sont en pointillés. Les longueurs obliques représentées ne sont pas à mesurer directement. Des vues de dessus et de face aident à dénombrer les cubes d’un empilement.\end{savoirfaire}
\begin{attention}Un dessin en perspective peut masquer des cubes. Un tableau de hauteurs, une vue ou une précision du texte doit permettre de décider ce qui est caché.\end{attention}
\clearpage
\phantomsection\label{cours-17-2}
\essSection{2}{Déplier un solide}
\begin{definition}Un patron est une figure plane qui, par pliage ou assemblage, donne la surface du solide sans chevauchement des faces. Deux faces qui se rejoignent ont des bords de même longueur.\end{definition}
\label{fig-cours-17-2-cours-patrons}\figureGeo{cours-patrons}

\begin{propriete}Un prisme droit de base triangulaire a deux triangles identiques et trois rectangles. Le patron d’un cylindre comporte deux disques de rayon $r$ et un rectangle de hauteur $h$ et de longueur $2\pi r$, le périmètre d’un disque.\end{propriete}
\begin{savoirfaire}[Vérifier avant de découper]Associer d’abord chaque face du patron à la perspective. Contrôler les longueurs des bords, puis imaginer le pliage. Pour un cube, six carrés seuls ne suffisent pas : certains assemblages se recouvrent au pliage.\end{savoirfaire}
\clearpage
\phantomsection\label{cours-17-3}
\essSection{3}{Volumes et unités}
\begin{propriete}$V_{\mathrm{cube}}=c^3$, $V_{\mathrm{pavé}}=L\ell h$, $V_{\mathrm{prisme}}=B h$, où $B$ est l’aire d’une base et $h$ la distance entre les deux bases. Les trois dimensions doivent être dans la même unité.\end{propriete}
\begin{center}\renewcommand{\arraystretch}{1.35}\begin{tabular}{c|ccc}\toprule
Unité & m$^3$ & dm$^3$ & cm$^3$\\\midrule
Vers la suivante & $\times1000$ & $\times1000$ & \\
\bottomrule\end{tabular}\end{center}
\begin{propriete}$1$ L $=1$ dm$^3$ ; $1$ mL $=1$ cm$^3$ ; $1$ m$^3=1000$ L. Les unités de capacité suivent les préfixes décimaux : $1$ L $=100$ cL $=1000$ mL.\end{propriete}
\begin{exemple}Un prisme de base triangulaire de base 6 cm et hauteur 4 cm a une aire de base de 12 cm$^2$. Pour une longueur de prisme de 10 cm, le volume vaut $12\times10=120$ cm$^3$.\end{exemple}
\begin{attention}Ne pas confondre la hauteur du triangle de base avec la hauteur du prisme. Une aire est en unités carrées, un volume en unités cubes.\end{attention}
\clearpage
\phantomsection\label{cours-17-4}
\essSection{4}{Disque et cylindre}
\begin{propriete}Un disque de rayon $r$ a pour aire $\pi r^2$. Le cylindre de rayon $r$ et de hauteur $h$ a pour volume $\pi r^2h$. Le rayon est la moitié du diamètre.\end{propriete}
\label{fig-cours-17-4-cours-cylindre}\figureGeo{cours-cylindre}

\begin{exemple}Pour $r=3$ cm et $h=8$ cm, l’aire de base est $9\pi$ cm$^2$ et le volume $72\pi$ cm$^3$, soit environ 226 cm$^3$ à l’unité.\end{exemple}
\begin{savoirfaire}[Comprendre la formule]On découpe le disque en secteurs très fins, que l’on dispose tête-bêche : on obtient presque un parallélogramme de base $\pi r$, la moitié du périmètre, et de hauteur $r$. Son aire vaut $\pi r\times r=\pi r^2$. Comme pour le prisme, le volume du cylindre est l’aire de base multipliée par la hauteur.\end{savoirfaire}
\begin{savoirfaire}[Exact puis approché]Conserver $\pi$ dans le résultat exact. Utiliser ensuite la calculatrice pour l’arrondi demandé. Ne pas remplacer trop tôt $\pi$ par 3,14, ce qui peut dégrader l’arrondi final.\end{savoirfaire}
\begin{attention}Pour peindre un objet, on calcule une aire de surface ; pour le remplir, un volume ou une capacité. Les données et la question indiquent la formule utile.\end{attention}

\clearpage\rubrique{Découverte}
\begin{blocLivre}{ex-17-1}\begin{exercice}[1]{Reconnaître un solide}
Identifier les solides A et B et donner le nombre de faces planes de chacun.
\label{fig-ex-17-1}\figureGeo{sol-1}
\end{exercice}\end{blocLivre}
\begin{blocLivre}{ex-17-2}\begin{exercice}[1]{Volume d’un cube}
Un cube a un côté de 4 cm. Calculer son volume.
\label{fig-ex-17-2}\figureGeo{sol-2}
\end{exercice}\end{blocLivre}
\begin{blocLivre}{ex-17-3}\begin{exercice}[1]{Volume d’un pavé}
Un pavé mesure 7 cm sur 5 cm sur 3 cm. Calculer son volume.
\label{fig-ex-17-3}\figureGeo{sol-3}
\end{exercice}\end{blocLivre}
\begin{blocLivre}{ex-17-4}\begin{exercice}[1]{Aire d’un disque}
Un disque a pour rayon 4 cm. Calculer son aire exacte.
\label{fig-ex-17-4}\figureGeo{sol-4}
\end{exercice}\end{blocLivre}
\begin{blocLivre}{ex-17-5}\begin{exercice}[1]{Volume et capacité}
Convertir 2,5 dm$^3$ en litres et en cm$^3$.
\label{fig-ex-17-5}\figureGeo{sol-5}
\end{exercice}\end{blocLivre}

\rubrique{Entraînement}
\begin{blocLivre}{ex-17-6}\begin{exercice}[2]{Prisme triangulaire}
La base triangulaire d’un prisme a une base de 6 cm et une hauteur de 4 cm. La longueur du prisme vaut 9 cm. Calculer son volume.
\label{fig-ex-17-6}\figureGeo{sol-6}
\end{exercice}\end{blocLivre}
\begin{blocLivre}{ex-17-7}\begin{exercice}[2]{Cylindre}
Un cylindre a un rayon de 2 cm et une hauteur de 10 cm. Calculer son volume exact, puis arrondir au cm$^3$.
\label{fig-ex-17-7}\figureGeo{sol-7}
\end{exercice}\end{blocLivre}
\begin{blocLivre}{ex-17-8}\begin{exercice}[2]{Diamètre donné}
Un cylindre a un diamètre de 6 cm et une hauteur de 5 cm. Calculer son volume exact.
\label{fig-ex-17-8}\figureGeo{sol-8}
\end{exercice}\end{blocLivre}
\begin{blocLivre}{ex-17-9}\begin{exercice}[2]{Un patron de prisme}
Un prisme droit a pour base un triangle de côtés 3, 4 et 5 cm et une longueur de 6 cm. Compléter son patron avec les dimensions des trois rectangles.
\label{fig-ex-17-9}\figureGeo{sol-9}
\end{exercice}\end{blocLivre}
\begin{blocLivre}{ex-17-10}\begin{exercice}[2]{Un patron de cylindre}
Un cylindre a un rayon de 3 cm et une hauteur de 7 cm. Donner les dimensions du rectangle de son patron.
\label{fig-ex-17-10}\figureGeo{sol-10}
\end{exercice}\end{blocLivre}

\rubrique{Problèmes}
\begin{blocLivre}{ex-17-11}\begin{exercice}[3]{Compter des cubes}
Le tableau de hauteurs vu de dessus indique trois colonnes au premier rang de hauteurs 2, 1, 3 ; au second rang, 1, 2, 0. Combien de cubes y a-t-il ?
\label{fig-ex-17-11}\figureGeo{sol-11}
\end{exercice}\end{blocLivre}
\begin{blocLivre}{ex-17-12}\begin{exercice}[3]{Remplir un aquarium}
Un aquarium intérieur mesure 60 cm sur 30 cm. On le remplit jusqu’à une hauteur de 25 cm. Quel volume d’eau faut-il en litres ?
\label{fig-ex-17-12}\figureGeo{sol-12}
\end{exercice}\end{blocLivre}
\begin{blocLivre}{ex-17-13}\begin{exercice}[3]{Comparer des récipients}
A est un pavé de 10 cm sur 8 cm sur 6 cm. B est un cylindre de rayon 4 cm et de hauteur 10 cm. Lequel contient le plus ?
\label{fig-ex-17-13}\figureGeo{sol-13}
\end{exercice}\end{blocLivre}
\begin{blocLivre}{ex-17-14}\begin{exercice}[3]{Compléter une perspective}
On a commencé la représentation en perspective cavalière d’un prisme droit ABCDEF dont la base ABC est un triangle rectangle en A : la face avant est en vraie grandeur et l’arête [BE] est tracée. Reproduire puis compléter le dessin, avec les arêtes cachées en pointillés. Citer ensuite les arêtes de même longueur que [BE] et la face parallèle à ABC.
\label{fig-ex-17-14}\figureGeo{sol-14}
\end{exercice}\end{blocLivre}
\begin{blocLivre}{ex-17-15}\begin{exercice}[3]{Une boîte ouverte}
Une boîte sans couvercle mesure 12 cm sur 8 cm sur 5 cm. Calculer sa capacité en mL et l’aire de carton utilisée, sans languettes.
\label{fig-ex-17-15}\figureGeo{sol-15}
\end{exercice}\end{blocLivre}

\clearpage
